\documentclass[10pt,a4paper]{amsart}
\usepackage[margin=19mm]{geometry}
\usepackage{amsmath,amssymb,mathtools}
\usepackage[scr=rsfs,cal=euler]{mathalfa}
\usepackage{xfrac}
\usepackage[unicode,hidelinks,bookmarks=false]{hyperref}
\newcommand{\ie}{i.e.\ }
\newcommand{\cf}{cf.\ }
\newcommand{\resp}{respectively\ }
\newcommand{\Subsection}{\subsection}
\providecommand{\nobreakdash}{\nobreak\hbox{-}}
\setlength{\emergencystretch}{2em}
\multlinegap0pt
\usepackage{cleveref}
\newtheorem{theorem}{Theorem}
\newtheorem{lemma}[theorem]{Lemma}
\newtheorem{proposition}[theorem]{Proposition}
\newcommand{\pr}{\mathbb{P}}
\newcommand{\ev}{\mathbb{E}}
\newcommand{\ent}{\mathcal{H}}
\newcommand{\cL}{\mathcal L}
\newcommand{\C}{\mathcal C}
\newcommand{\X}{\mathcal X}
\newcommand{\logt}{\log_2}
\DeclarePairedDelimiter{\abs}{\lvert}{\rvert}
\title[Cycle-factors of regular graphs via entropy: the expected number of cycles (Proposition 3.4)]{Cycle-factors of regular graphs via entropy: the expected number of cycles of a random cycle-factor (Proposition 3.4)}
\author{Micha Christoph, Nemanja Dragani\'c, Ant\'onio Gir\~ao, Eoin Hurley, Lukas Michel, Alp M\"uyesser}
\date{Source: arXiv:2507.19417v3 (CC BY 4.0)}
\begin{document}
\maketitle

\noindent\emph{Source and licence.} Cycle-factors of regular graphs via entropy. Version arXiv:2507.19417v3 (CC BY 4.0), distributed under the Creative Commons Attribution 4.0 International licence, as recorded in the arXiv licence metadata for this exact version and shown on the abstract page https://arxiv.org/abs/2507.19417v3. Full rights notice: licenses/arxiv-2507.19417-CC-BY-4.0.txt.\par\smallskip

\noindent\textit{Editorial scope.} Counted unit: the proposition printed as Proposition 3.4 of the article (source0.tex lines 246--251, source label \texttt{prop:boundexpectednumberofcycles}), delivered together with the article's own complete original proof of it (source0.tex lines 253--315, one \texttt{proof} environment of 63 lines that proves the stated upper bound on the expected number of cycles with no gap). No mathematics is written, repaired or re-derived: statement and proof are the article's own text line for line, in the article's own notation ($C$, $G$, $d$, $n$, $\sigma$, $\C$, $\ent$, $a(i,\sigma',\tau)$, $\ell(i,\sigma',\tau)$, $\cL$). Required context retained uncounted: the article's Lemma 3.3, the entropy-loss bound that the counted proof invokes, in statement form (lines 217--222, source label \texttt{lem:entlossskew}). Every definition, display and label of that line range is retained unchanged. Omitted from this package and not redistributed: the article's preamble and title page with its six-author block and affiliations (lines 1--112); its abstract and the rest of its front matter (lines 113--116); its introduction with the announcements of its two main theorems, its two bipartite-graph inequalities and its algorithmic corollaries (lines 117--216); its entropy section apart from the retained lemma (lines 223--245); the paragraph introducing the counted proposition (line 244); the article's further consequences of the counted proposition, namely the proofs of its main theorems, its one-sided concentration bound, the entropy-function section and its open problems (lines 316--382). Nothing inside a retained excerpt is omitted or altered: every retained body is delivered exactly as the article's lines are written, so no omission receipt of any kind accompanies this package. Citations: the retained excerpts carry no citation command at all, so no bibliography entry of the article is transcribed and no reference is redistributed. Reference adaptation: none was needed and none was made. Every reference inside the delivered unit resolves inside it: the label lem:entlossskew of the retained lemma, which the counted proof cites with \texttt{\textbackslash cref}, and the labels of the counted statement itself. No unresolved reference is produced. Printed slips kept literal: no printed word, symbol, hypothesis or number of the counted statement or of its proof was altered, and the retained text keeps the article's own \texttt{\textbackslash qedhere} at the end of its final display. Layout and class: the article's own class is the standard LaTeX \texttt{article} class with the option 11pt, using mathpazo and loading among others amsmath, amsthm, mathtools, thmtools, bbm, cleveref and hyperref; no publisher class or style file is used or shipped, so none travels with this package and none is redistributed. The wrapper is the lane's pinned \texttt{amsart} editorial wrapper at 10pt on A4, which adds the article's own macro definitions that the retained text needs ($\pr$, $\ev$, $\ent$, $\cL$, $\C$, $\X$, $\logt$ and the paired delimiter \texttt{\textbackslash abs}) and loads the \texttt{cleveref} package that the retained proof needs for its \texttt{\textbackslash cref}. The delivered numbering is the unit's own: the retained Lemma 3.3 prints as Lemma 1 and the counted Proposition 3.4 prints as Proposition 2. The article's 11pt measure and its Palatino text font differ from the wrapper's 10pt Latin Modern; those are page-appearance differences of the wrapper only, and no formula, symbol, hypothesis or argument changes. Assets: no retained span contains a figure environment, a graphics-inclusion command, a PDF-page inclusion, a table or a TikZ picture, and the article has none anywhere in its source; no image file is copied into this package, the package declares no asset dependency and no image file is redistributed with it. Licence: version arXiv:2507.19417v3 of this article is distributed under the Creative Commons Attribution 4.0 International licence, as recorded in the arXiv licence metadata for this exact version and shown on the abstract page https://arxiv.org/abs/2507.19417v3. A full rights notice is delivered at licenses/arxiv-2507.19417-CC-BY-4.0.txt. Characters: every retained excerpt is pure ASCII, so the delivered engine, XeLaTeX with its default Latin Modern text font, reports no missing character; the article's own accented letters all lie in its author block and abstract, which are not retained. No glyph is substituted and no word is rewritten.
\begin{lemma}\label{lem:entlossskew}
    Let $X$ be a random variable taking values in a finite set $\X$, and let $\ell \coloneqq \logt \abs{\X} - \ent(X)$. Then, for all $x \in \X$ we have
    \[
        \pr(X = x) \le \frac{2}{\abs{\X}} + \ell.
    \]
\end{lemma}

\begin{proposition}\label{prop:boundexpectednumberofcycles}
    Let $C$ be a random cycle-factor of a directed $d$-regular graph $G$ on $n$ vertices. Then, the expected number of cycles of $C$ is at most
    \[
        2 \cdot \frac{n}{d} (\log d + 1) + \frac{n}{d} \logt(d!) - \ent(C).
    \]
\end{proposition}

\begin{proof}
    Identify the vertices of $G$ with $[n]$. We will view $C$ as the random permutation $\sigma$ of $[n]$ where $\sigma(i)$ is the out-neighbour of $i$ in $C$ for all $i \in [n]$. In particular, the number of cycles of $C$ is exactly $\abs{\sigma}$, and $\ent(C) = \ent(\sigma)$. Also, let $\C$ be the collection of all permutations $\sigma'$ of $[n]$ with $(i,\sigma'(i)) \in E(G)$ for all $i \in [n]$, so $\C$ corresponds to the collection of all cycle-factors of $G$.

    For every permutation $\tau \in S_n$, the chain rule for entropy implies that
    \[
        \ent(\sigma) = \sum_{j=1}^n \ent(\sigma(\tau(j)) \mid \sigma(\tau(1)), \dots, \sigma(\tau(j-1))).
    \]
    For $i \in [n]$, let $\tau_{<i} \coloneqq \tau(1), \dots, \tau(j-1)$, where $j \coloneqq \tau^{-1}(i)$. We will use the notation $\sigma(\tau_{<i}) \coloneqq \sigma(\tau(1)), \dots, \sigma(\tau(j-1))$. If we reorder the above sum, we get
    \[
        \ent(\sigma) = \sum_{i=1}^n \ent(\sigma(i) \mid \sigma(\tau_{<i})).
    \]
    Expanding the definition of conditional entropy, this yields
    \[
        \ent(\sigma) = \sum_{i=1}^n \sum_\pi \pr(\sigma(\tau_{<i}) = \pi) \cdot \ent(\sigma(i) \mid \sigma(\tau_{<i}) = \pi),
    \]
    where $\pi$ ranges over all possible values of $\sigma(\tau_{<i})$. Observe that $\pr(\sigma(\tau_{<i}) = \pi)$ is exactly the sum of the probabilities $\pr(\sigma = \sigma')$ over all $\sigma' \in \C$ with $\sigma'(\tau_{<i}) = \pi$. Therefore,
    \[
        \ent(\sigma) = \sum_{i=1}^n \sum_{\sigma' \in \C} \pr(\sigma = \sigma') \cdot \ent(\sigma(i) \mid \sigma(\tau_{<i}) = \sigma'(\tau_{<i})).
    \]
    For all $i \in [n]$, $\sigma' \in \C$, and $\tau \in S_n$, let $a(i,\sigma',\tau) \coloneqq \abs{N_G^+(i) \setminus \sigma'(\tau_{<i})}$ be the number of \emph{available out-neighbours} of $i$ in $G$, that is, the number of out-neighbours of $i$ in $G$ that are not contained in $\sigma'(\tau_{<i})$. If $\sigma(\tau_{<i}) = \sigma'(\tau_{<i})$, then $\sigma(i)$ must be one of these available out-neighbours, and so $\ent(\sigma(i) \mid \sigma(\tau_{<i}) = \sigma'(\tau_{<i})) \le \logt a(i,\sigma',\tau)$. Denote the loss in entropy in this inequality by $\ell(i,\sigma',\tau) \coloneqq \logt a(i,\sigma',\tau) - \ent(\sigma(i) \mid \sigma(\tau_{<i}) = \sigma'(\tau_{<i}))$. Then,
    \[
        \ent(\sigma) = \sum_{i=1}^n \sum_{\sigma' \in \C} \pr(\sigma = \sigma') \cdot (\logt a(i,\sigma',\tau) - \ell(i,\sigma',\tau)).
    \]
    Since this holds for all $\tau \in S_n$, we may pick $\tau \in S_n$ uniformly at random. This yields
    \begin{align*}
        \ent(\sigma) & = \ev_\tau\left(\sum_{i=1}^n \sum_{\sigma' \in \C} \pr(\sigma = \sigma') \cdot (\logt a(i,\sigma',\tau) - \ell(i,\sigma',\tau))\right) \\
        & = \sum_{i=1}^n \sum_{\sigma' \in \C} \pr(\sigma = \sigma') \cdot \ev_\tau(\logt a(i,\sigma',\tau)) - \cL
    \end{align*}
    where
    \[
        \cL \coloneqq \sum_{i=1}^n \sum_{\sigma' \in \C} \pr(\sigma = \sigma') \cdot \ev_\tau(\ell(i,\sigma',\tau))
    \]
    denotes the loss in entropy compared to a uniform distribution.
    
    For a fixed $i$ and $\sigma'$, note that $a(i,\sigma',\tau)$ is the number of out-neighbours $j \in N_G^+(i)$ of $i$ in $G$ such that $(\sigma')^{-1}(j)$ does not appear before $i$ in $\tau$. Since $\tau$ is a uniformly random permutation, this number is distributed uniformly at random in $[d]$, and so it follows that
    \[
        \ent(\sigma) = \sum_{i=1}^n \sum_{\sigma' \in \C} \pr(\sigma = \sigma') \cdot \left(\frac{1}{d} \sum_{t=1}^d \logt t\right) - \cL = \frac{n}{d} \logt(d!) - \cL.
    \]
    In particular, we get that $\cL = n \logt(d!) / d - \ent(\sigma) = n \logt(d!) / d - \ent(C)$.

    We now bound the number of cycles $\abs{\sigma}$ of $\sigma$ using a very similar calculation. For every permutation $\tau \in S_n$, we say that $\sigma(i)$ closes a cycle if $\tau_{<i}$ contains all other vertices that are on the same cycle as $i$ in $\sigma$. That is, adding $\sigma(i)$ to $\sigma(\tau_{<i})$ creates one additional cycle. Note that $\abs{\sigma}$ is equal to the number of vertices $i \in [n]$ for which $\sigma(i)$ closes a cycle, and so
    \begin{align*}
        \ev_\sigma(\abs{\sigma}) & = \sum_{i = 1}^n \pr(\sigma(i) \text{ closes a cycle}) \\
        & = \sum_{i=1}^n \sum_\pi \pr(\sigma(\tau_{<i}) = \pi) \cdot \pr(\sigma(i) \text{ closes a cycle} \mid \sigma(\tau_{<i}) = \pi) \\
        & = \sum_{i=1}^n \sum_{\sigma' \in \C} \pr(\sigma = \sigma') \cdot \pr(\sigma(i) \text{ closes a cycle} \mid \sigma(\tau_{<i}) = \sigma'(\tau_{<i})).
    \end{align*}
    If $\sigma(\tau_{<i}) = \sigma'(\tau_{<i})$, either $\sigma(i)$ cannot close a cycle, or $\sigma(i)$ needs to be one specific available out-neighbour $j \in N_G^+(i) \setminus \sigma'(\tau_{<i})$ in order to close a cycle. So, by \cref{lem:entlossskew} we know that
    \begin{align*}
        \pr(\sigma(i) \text{ closes a cycle} \mid \sigma(\tau_{<i}) = \sigma'(\tau_{<i})) & \le \max_{j \in N_G^+(i)} \pr(\sigma(i) = j \mid \sigma(\tau_{<i}) = \sigma'(\tau_{<i})) \\
        & \le \frac{2}{a(i,\sigma',\tau)} + \ell(i,\sigma',\tau).
    \end{align*}
    This implies that
    \[
        \ev_\sigma(\abs{\sigma}) \le \sum_{i=1}^n \sum_{\sigma' \in \C} \pr(\sigma = \sigma') \cdot \left(\frac{2}{a(i,\sigma',\tau)} + \ell(i,\sigma',\tau)\right).
    \]
    Since this holds for all $\tau \in S_n$, we may again pick $\tau \in S_n$ uniformly at random. As noted before, $a(i,\sigma',\tau)$ is then distributed uniformly at random in $[d]$, and so
    \begin{align*}
        \ev_\sigma(\abs{\sigma}) & \le \ev_\tau\left(\sum_{i=1}^n \sum_{\sigma' \in \C} \pr(\sigma = \sigma') \cdot \left(\frac{2}{a(i,\sigma',\tau)} + \ell(i,\sigma',\tau)\right)\right) \\
        & = \sum_{i=1}^n \sum_{\sigma' \in \C} \pr(\sigma = \sigma') \cdot \ev_\tau\left(\frac{2}{a(i,\sigma',\tau)}\right) + \cL \\
        & = \sum_{i=1}^n \sum_{\sigma' \in \C} \pr(\sigma = \sigma') \cdot \left(\frac{1}{d} \sum_{t = 1}^d \frac{2}{t}\right) + \cL \\
        & \le 2 \cdot \frac{n}{d} (\log d + 1) + \cL = 2 \cdot \frac{n}{d} (\log d + 1) + \frac{n}{d} \logt(d!) - \ent(C). \qedhere
    \end{align*}
\end{proof}

\end{document}
